\documentclass[UTF8,fontset=none,zihao=5,twoside,openany]{ctexbook}
\usepackage{bookstyle}
\usepackage{float}
% 样张中让所有图固定在原位，便于前后对照
\let\oldfigure\figure \let\endoldfigure\endfigure
\renewenvironment{figure}[1][]{\oldfigure[H]}{\endoldfigure}
\setcounter{chapter}{4}
\pagestyle{plain}
\begin{document}

\section*{样图 1：图 4-9 预填充与解码}
\textbf{问题：}“KV 缓存”标注贴在长条左端；右侧红色注释外挂，使整幅图比正文栏偏左约 6 mm。\\
\textbf{改法：}整幅图以正文栏中线为轴对称排布；“KV 缓存”标注居中并说明其增长方式；“计算受限 / 访存受限”并入下方括号；外挂注释改为底部居中的图例；首个 token 改由预填充产生，与 TTFT 的定义一致。

\noindent\textbf{修改前}
\begin{figure}[htbp]
\centering
\begin{tikzpicture}[fig, x=1mm, y=1mm]
  \tikzset{tok/.style={draw=figline, fill=white, minimum size=4.2mm, inner sep=0pt, font=\figfont\footnotesize},
           newtok/.style={draw=accent, fill=accentpale, minimum size=4.2mm, inner sep=0pt, font=\figfont\footnotesize}}
  % ---------- 计算行：整体以 x=0 为中心，范围约 -55 ~ 55 ----------
  \node[blk, minimum width=38mm, minimum height=9mm] (pf) at (-36,0) {预填充\\[-1pt]\footnotesize 并行处理全部输入};
  \foreach \i in {1,...,5} {
    \node[tok] (x\i) at (-57+\i*5,14) {x\textsubscript{\i}};
    \draw[->, line width=0.4pt] (x\i.south) -- (x\i.south |- pf.north);
  }
  \node[newtok] (y1) at (-21,14) {y\textsubscript{1}};
  \draw[->, line width=0.4pt] (y1.south |- pf.north) -- (y1.south);
  \foreach \k [evaluate=\k as \n using int(\k+1)] in {1,...,5} {
    \pgfmathsetmacro{\xx}{-6+(\k-1)*12}
    \node[blk, minimum width=9mm, minimum height=9mm, inner sep=1pt, font=\figfont\scriptsize] (dc\k) at (\xx,0) {解码};
    \node[newtok] (y\n) at (\xx,14) {y\textsubscript{\n}};
    \draw[->, line width=0.4pt] (dc\k.north) -- (y\n.south);
  }
  \draw[->, line width=0.4pt] (pf.east) -- (dc1.west);
  \foreach \k [evaluate=\k as \n using int(\k+1)] in {1,...,4} \draw[->, line width=0.4pt] (dc\k.east) -- (dc\n.west);
  \node[font=\figfont\small] at (51.5,0) {…};
  \node[note, anchor=south] at (-36,17.5) {输入提示（T 个 token）};
  \node[note, anchor=south] at (18,17.5) {逐个生成输出 token};
  % ---------- KV 缓存：居中标注，分段表示逐步增长 ----------
  \node[blkw, minimum width=110mm, minimum height=7mm] (kv) at (0,-16) {KV 缓存（提示部分一次写入，生成时每步追加一项）};
  \draw[->, accent, line width=0.7pt] (pf.south) -- (pf.south |- kv.north);
  \foreach \k in {1,...,5} \draw[<->, warn, line width=0.7pt] (dc\k.south) -- (dc\k.south |- kv.north);
  % ---------- 时间轴与指标 ----------
  \draw[->, line width=0.5pt, rulegray] (-55,-25) -- (57,-25) node[right, note, text=ink] {时间};
  \draw[decorate, decoration={brace, mirror, amplitude=3pt}, line width=0.6pt] (-55,-27) -- (-17,-27)
    node[midway, below=3pt, note] {TTFT：首 token 延迟\\[-1pt]（计算受限）};
  \draw[decorate, decoration={brace, mirror, amplitude=3pt}, line width=0.6pt] (-10.5,-27) -- (-1.5,-27)
    node[midway, below=3pt, note] {TPOT\\[-1pt]（访存受限）};
  % ---------- 居中图例 ----------
  \begin{scope}[shift={(0,-44)}]
    \draw[->, accent, line width=0.7pt] (-44,0) -- (-38,0); \node[note, anchor=west] at (-37,0) {一次性写入提示的 K/V};
    \draw[<->, warn, line width=0.7pt] (2,0) -- (8,0);      \node[note, anchor=west] at (9,0) {每步读取历史 K/V 并追加新 K/V};
  \end{scope}
\end{tikzpicture}
\caption{一次推理请求的两个阶段。预填充并行处理全部输入、写入 KV 缓存并给出首个 token；之后每个解码步只生成一个 token，却要读取全部权重与已有缓存}
\label{fig:prefill}
\end{figure}

\Needspace{80mm}\noindent\textbf{修改后}
\begin{figure}[htbp]
\centering
\begin{tikzpicture}[fig, x=1mm, y=1mm]
  \tikzset{tok/.style={draw=figline, fill=white, minimum size=4.2mm, inner sep=0pt, font=\figfont\scriptsize},
           newtok/.style={draw=accent, fill=accentpale, minimum size=4.2mm, inner sep=0pt, font=\figfont\scriptsize}}
  % ---------- 计算行：整体以 x=0 为中心，范围约 -55 ~ 55 ----------
  \node[blk, minimum width=38mm, minimum height=9mm] (pf) at (-36,0) {预填充\\[-1pt]\footnotesize 并行处理全部输入};
  \foreach \i in {1,...,5} {
    \node[tok] (x\i) at (-57+\i*5,14) {$x_{\i}$};
    \draw[->, line width=0.4pt] (x\i.south) -- (x\i.south |- pf.north);
  }
  \node[newtok] (y1) at (-21,14) {$y_{1}$};
  \draw[->, line width=0.4pt] (y1.south |- pf.north) -- (y1.south);
  \foreach \k [evaluate=\k as \n using int(\k+1)] in {1,...,5} {
    \pgfmathsetmacro{\xx}{-6+(\k-1)*12}
    \node[blk, minimum width=9mm, minimum height=9mm, inner sep=1pt, font=\figfont\scriptsize] (dc\k) at (\xx,0) {解码};
    \node[newtok] (y\n) at (\xx,14) {$y_{\n}$};
    \draw[->, line width=0.4pt] (dc\k.north) -- (y\n.south);
  }
  \draw[->, line width=0.4pt] (pf.east) -- (dc1.west);
  \foreach \k [evaluate=\k as \n using int(\k+1)] in {1,...,4} \draw[->, line width=0.4pt] (dc\k.east) -- (dc\n.west);
  \node[font=\figfont\small] at (51.5,0) {$\cdots$};
  \node[note, anchor=south] at (-36,17.5) {输入提示（$T$ 个 token）};
  \node[note, anchor=south] at (18,17.5) {逐个生成输出 token};
  % ---------- KV 缓存：居中标注，分段表示逐步增长 ----------
  \node[blkw, minimum width=110mm, minimum height=7mm] (kv) at (0,-16) {KV 缓存（提示部分一次写入，生成时每步追加一项）};
  \draw[->, accent, line width=0.7pt] (pf.south) -- (pf.south |- kv.north);
  \foreach \k in {1,...,5} \draw[<->, warn, line width=0.7pt] (dc\k.south) -- (dc\k.south |- kv.north);
  % ---------- 时间轴与指标 ----------
  \draw[->, line width=0.5pt, rulegray] (-55,-25) -- (57,-25) node[right, note, text=ink] {时间};
  \draw[decorate, decoration={brace, mirror, amplitude=3pt}, line width=0.6pt] (-55,-27) -- (-17,-27)
    node[midway, below=3pt, note] {TTFT：首 token 延迟\\[-1pt]\textcolor{rulegray}{（计算受限）}};
  \draw[decorate, decoration={brace, mirror, amplitude=3pt}, line width=0.6pt] (-10.5,-27) -- (-1.5,-27)
    node[midway, below=3pt, note] {TPOT\\[-1pt]\textcolor{rulegray}{（访存受限）}};
  % ---------- 居中图例 ----------
  \begin{scope}[shift={(0,-44)}]
    \draw[->, accent, line width=0.7pt] (-44,0) -- (-38,0); \node[note, anchor=west] at (-37,0) {一次性写入提示的 K/V};
    \draw[<->, warn, line width=0.7pt] (2,0) -- (8,0);      \node[note, anchor=west] at (9,0) {每步读取全部权重与缓存};
  \end{scope}
\end{tikzpicture}
\caption{（修改稿）一次推理请求的两个阶段。预填充并行处理全部输入、写入 KV 缓存并给出首个 token；之后每个解码步只生成一个 token，却要读取全部权重与已有缓存}
\label{fig:s-prefill}
\end{figure}


\clearpage
\section*{样图 2：图 4-12 综合案例的设计流程}
\textbf{问题：}主流程靠左，取值文字外挂在右侧且不对齐，底部说明偏左，整体重心偏移。\\
\textbf{改法：}改成左右两栏的对照表结构（设计环节 / 本案例中的取值），加栏头；取值放进细边框卡片并与环节逐行对齐；图例居中。

\noindent\textbf{修改前}
\begin{figure}[htbp]
\centering
\begin{tikzpicture}[fig, x=1mm, y=1mm]
  % 两栏对齐：左栏为设计环节，右栏为本案例取值；整体以 x=0 为中心
  \node[note, font=\figfont\footnotesize\bfseries] at (-16,9) {设计环节};
  \node[note, font=\figfont\footnotesize\bfseries] at (32,9) {本案例中的取值};
  \foreach \i/\t/\c in {
      0/业务需求/{200 名学生，高峰 30 路并发\\回答注明出处，数据不外传},
      1/服务指标/{P95 TTFT ≤ 2 s，输出 ≥ 15 token/s\\引用可核对},
      2/模型与推理引擎/{RAG 工作流 + 8B 开源模型\\vLLM 类推理引擎},
      3/显存与带宽估算/{4 位权重约 5 GB\\KV 缓存每请求约 0.5 GiB},
      4/硬件与网络/{单张 24 GB 显卡\\显存带宽约 1 TB/s}} {
    \ifnum\i=3
      \node[blka, minimum width=32mm, minimum height=10mm] (s\i) at (-16,-\i*14) {\t};
    \else
      \node[blk, minimum width=32mm, minimum height=10mm] (s\i) at (-16,-\i*14) {\t};
    \fi
    \node[blkw, draw=rulegray, minimum width=50mm, minimum height=10mm, font=\figfont\footnotesize, align=left, text width=46mm] (v\i) at (32,-\i*14) {\c};
    \draw[densely dotted, rulegray, line width=0.6pt] (s\i.east) -- (v\i.west);
  }
  \foreach \i [evaluate=\i as \j using int(\i+1)] in {0,...,3}
    \draw[flow] (s\i.south) -- (s\j.north);
  % 实测反馈回路（左侧）
  \node[blkw, minimum width=15mm, minimum height=10mm, font=\figfont\footnotesize, align=center] (m) at (-50,-35) {实测\\验证};
  \draw[->, dashed, line width=0.6pt] (s4.west) -| (m.south);
  \draw[->, dashed, line width=0.6pt] (m.north) |- (s1.west);
  % 居中图例
  \begin{scope}[shift={(0,-67)}]
    \draw[flow] (-40,0) -- (-34,0); \node[note, anchor=west, font=\figfont\scriptsize] at (-33,0) {设计方向：自上而下};
    \draw[->, dashed, line width=0.6pt] (2,0) -- (8,0); \node[note, anchor=west, font=\figfont\scriptsize] at (9,0) {用测量结果修正指标与方案};
  \end{scope}
\end{tikzpicture}
\caption{自上而下的系统设计流程及其在课程知识问答助手中的取值。高亮的估算环节决定了方案能否成立}
\label{fig:case-flow}
\end{figure}

\Needspace{80mm}\noindent\textbf{修改后}
\begin{figure}[htbp]
\centering
\begin{tikzpicture}[fig, x=1mm, y=1mm]
  % 两栏对齐：左栏为设计环节，右栏为本案例取值；整体以 x=0 为中心
  \node[note, font=\figfont\footnotesize\bfseries, text=rulegray] at (-16,9) {设计环节};
  \node[note, font=\figfont\footnotesize\bfseries, text=rulegray] at (32,9) {本案例中的取值};
  \foreach \i/\t/\c in {
      0/业务需求/{200 名学生，高峰 30 路并发\\回答注明出处，数据不外传},
      1/服务指标/{P95 TTFT $\le$ 2 s，输出 $\ge$ 15 token/s\\引用可核对},
      2/模型与推理引擎/{RAG 工作流 + 8B 开源模型\\vLLM 类推理引擎},
      3/显存与带宽估算/{4 位权重约 5 GB\\KV 缓存每请求约 0.5 GiB},
      4/硬件与网络/{单张 24 GB 显卡\\显存带宽约 1 TB/s}} {
    \ifnum\i=3
      \node[blka, minimum width=32mm, minimum height=10mm] (s\i) at (-16,-\i*14) {\t};
    \else
      \node[blk, minimum width=32mm, minimum height=10mm] (s\i) at (-16,-\i*14) {\t};
    \fi
    \node[blkw, draw=rulegray, minimum width=50mm, minimum height=10mm, font=\figfont\footnotesize, align=left, text width=46mm] (v\i) at (32,-\i*14) {\c};
    \draw[densely dotted, rulegray, line width=0.6pt] (s\i.east) -- (v\i.west);
  }
  \foreach \i [evaluate=\i as \j using int(\i+1)] in {0,...,3}
    \draw[flow] (s\i.south) -- (s\j.north);
  % 实测反馈回路（左侧）
  \node[blkw, minimum width=15mm, minimum height=10mm, font=\figfont\footnotesize, align=center] (m) at (-50,-35) {实测\\验证};
  \draw[->, dashed, line width=0.6pt] (s4.west) -| (m.south);
  \draw[->, dashed, line width=0.6pt] (m.north) |- (s1.west);
  % 居中图例
  \begin{scope}[shift={(0,-67)}]
    \draw[flow] (-40,0) -- (-34,0); \node[note, anchor=west, font=\figfont\scriptsize] at (-33,0) {设计方向：自上而下};
    \draw[->, dashed, line width=0.6pt] (2,0) -- (8,0); \node[note, anchor=west, font=\figfont\scriptsize] at (9,0) {用测量结果修正指标与方案};
  \end{scope}
\end{tikzpicture}
\caption{（修改稿）自上而下的系统设计流程及其在课程知识问答助手中的取值。高亮的估算环节决定了方案能否成立}
\label{fig:s-case}
\end{figure}


\clearpage
\section*{样图 3：图 4-11 应用分层与五类工具}
\textbf{问题：}右侧散落的工具名与左侧图形不对称，“向量数据库”“部署工具”漂浮在箭头旁，读者难以对应。\\
\textbf{改法：}微调并入同一行，全图等宽居中；五类工具改为贴在所属层右上角的“标签”，与层一一对应，不再外挂。

\noindent\textbf{修改前}
% 图 4-6 第二稿：遵循 nature-figure 矢量示意图规范
%  - 一种字族、两种字号：层名 \small 加粗（唯一加粗处），说明与工具 \footnotesize 常规
%  - 全部文字同一深色；唯一强调色只标“工具类别”这一本图要传达的信息
%  - 同一角色的框等宽、上下边对齐；每个框内部顺序相同：层名 → 一行说明 → （底部）工具类别
%  - 连线标签与其他文字保持间距；虚线只表示一种含义：离线
\begin{figure}[htbp]
\centering
\begin{tikzpicture}[fig, x=1mm, y=1mm]
  % \layerbox{x0}{x1}{ytop}{h}{填充}{边框样式}{层名}{说明}{工具类别，可空}
  \newcommand{\layerbox}[9]{%
    \begin{scope}
      \path[clip, rounded corners=1.5pt] (#1,#3) rectangle (#2,#3-#4);
      \fill[#5] (#1,#3) rectangle (#2,#3-#4);
      \if\relax\detokenize{#9}\relax\else
        \fill[accentpale] (#1,#3-#4) rectangle (#2,#3-#4+5);
        \draw[accent, line width=0.4pt] (#1,#3-#4+5) -- (#2,#3-#4+5);
      \fi
    \end{scope}
    \draw[figline, rounded corners=1.5pt, #6] (#1,#3) rectangle (#2,#3-#4);
    \if\relax\detokenize{#9}\relax
      \node[align=center, text=ink] at ({(#1+#2)/2},{#3-#4/2})
        {{\small\bfseries #7}\\[0.5pt]{\footnotesize #8}};
    \else
      \node[align=center, text=ink] at ({(#1+#2)/2},{#3-(#4-5)/2})
        {{\small\bfseries #7}\\[0.5pt]{\footnotesize #8}};
      \node[text=ink, font=\figfont\footnotesize] at ({(#1+#2)/2},{#3-#4+2.5}) {工具：#9};
    \fi
  }
  % ---- 第 1 层 ----
  \layerbox{-55}{55}{0}{11}{figfill}{}{接口层}{网页、移动端、开放 API}{}
  % ---- 第 2 层 ----
  \layerbox{-55}{55}{-19}{16}{figfill}{}{编排层}{工作流与智能体、提示词组织、权限校验、安全审查、日志}{编排平台、应用开发框架}
  % ---- 第 3 层：三个等宽框，等间距 ----
  \layerbox{-55}{-23}{-48}{16}{figfill}{}{知识层}{向量检索、搜索、文档库}{向量数据库}
  \layerbox{-16}{16}{-48}{16}{figfill}{}{模型服务层}{推理引擎、模型 API}{部署工具}
  \layerbox{23}{55}{-48}{16}{white}{dashed}{微调}{离线训练，产出新权重}{微调工具}
  % ---- 第 4 层 ----
  \layerbox{-55}{55}{-76}{11}{white}{}{基础设施层}{计算芯片、存储、互联、驱动与软件栈（第\,4.2\,节）}{}
  % ---- 连线：主路径实线，离线虚线 ----
  \draw[<->, line width=0.6pt] (0,-11) -- (0,-19);
  \draw[<->, line width=0.6pt] (-39,-35) -- (-39,-48) node[midway, right=1.5mm, note] {检索资料};
  \draw[<->, line width=0.6pt] (0,-35)   -- (0,-48)   node[midway, right=1.5mm, note] {生成回答};
  \draw[->, dashed, line width=0.6pt] (23,-56) -- (16,-56);
  \foreach \x in {-39,0,39} \draw[->, line width=0.6pt] (\x,-64) -- (\x,-76);
\end{tikzpicture}
\caption{大模型应用的典型分层与五类工具的位置。每个框底部的色带标出该层对应的工具类别。编排层发起检索与模型调用；微调（虚线框）在线下进行，产出的权重交给模型服务层加载}
\label{fig:app-arch}
\end{figure}

\Needspace{80mm}\noindent\textbf{修改后}
\begin{figure}[htbp]
\centering
\begin{tikzpicture}[fig, x=1mm, y=1mm]
  \tikzset{chip/.style={draw=accent, fill=white, text=accent, rounded corners=1pt, inner sep=1.2pt,
                        font=\figfont\scriptsize, anchor=south east}}
  % 全部层宽 104 mm，以 x=0 为中心；工具类别以“标签”贴在所属层的右上角
  \node[blkw, minimum width=104mm, minimum height=9mm] (ui) at (0,0) {{\bfseries 接口层}\quad 网页 · 移动端 · 开放 API};
  \node[blka, minimum width=104mm, minimum height=11mm] (orch) at (0,-19) {{\bfseries 编排层}\\[-1pt]\footnotesize 工作流 / 智能体 · 提示词组织 · 权限校验 · 安全审查 · 日志};
  \node[blk, minimum width=36mm, minimum height=11mm] (know) at (-34,-40) {{\bfseries 知识层}\\[-1pt]\footnotesize 向量数据库 · 搜索 · 文档库};
  \node[blk, minimum width=36mm, minimum height=11mm] (model) at (6,-40) {{\bfseries 模型服务层}\\[-1pt]\footnotesize 推理引擎 · 模型 API};
  \node[blkw, dashed, minimum width=22mm, minimum height=11mm] (ft) at (41,-40) {{\bfseries 微调}\\[-1pt]\footnotesize 离线进行};
  \node[blkw, minimum width=104mm, minimum height=9mm] (infra) at (0,-59) {{\bfseries 基础设施层}\quad 计算芯片 · 存储 · 互联 · 软件栈};
  % 调用关系
  \draw[<->, line width=0.6pt] (ui.south) -- (orch.north);
  \draw[<->, line width=0.6pt] (orch.south -| know.north) -- (know.north) node[midway, right, note, font=\figfont\scriptsize] {检索资料};
  \draw[<->, line width=0.6pt] (orch.south -| model.north) -- (model.north) node[midway, right, note, font=\figfont\scriptsize] {生成回答};
  \draw[->, dashed, line width=0.6pt] (ft.west) -- (model.east) node[midway, above, note, font=\figfont\scriptsize] {权重};
  \draw[->, line width=0.6pt] (know.south) -- (know.south |- infra.north);
  \draw[->, line width=0.6pt] (model.south) -- (model.south |- infra.north);
  % 工具类别标签
  \node[chip] at ([xshift=-1.5mm]orch.north east) {编排平台 · 应用开发框架};
  \node[chip] at ([xshift=-1.5mm]know.north east) {向量数据库};
  \node[chip, anchor=north east] at ([xshift=-1.5mm]model.south east) {部署工具};
  \node[chip] at ([xshift=-1.5mm]ft.north east) {微调工具};
  \node[chip, draw=rulegray, text=rulegray] at ([xshift=-1.5mm]infra.north east) {见第 4.2 节};
  % 居中图例
  \node[chip, anchor=center] at (-20,-69) {标签};
  \node[note, anchor=west, font=\figfont\scriptsize] at (-16,-69) {五类工具各自所处的层};
\end{tikzpicture}
\caption{（修改稿）大模型应用的典型分层与五类工具的位置。编排层居中调度：检索资料和调用模型都是它发起的调用；微调在线下进行，产出的权重交给模型服务层加载}
\label{fig:s-app}
\end{figure}


\IfFileExists{figs/ai/prefill_decode.png}{%
\clearpage
\section*{样图 4：AI 生成插图（GPT Image，经 Codex）}
同一主题用图像生成模型绘制，仅供与 TikZ 版本比较风格与文字可读性。
\begin{figure}[H]\centering
\includegraphics[width=\linewidth]{figs/ai/prefill_decode.png}
\caption{（AI 生成）预填充与解码}
\end{figure}
\IfFileExists{figs/ai/request_journey.png}{%
\begin{figure}[H]\centering
\includegraphics[width=\linewidth]{figs/ai/request_journey.png}
\caption{（AI 生成）一次大模型请求的旅程}
\end{figure}}{}}{}

\end{document}
